% ECONOMICS_PROBLEM: {"id": "H11-328", "title": "Building effective state capacity", "jel_code": "H11", "source_id": "OP-0032"}
% FIRST_POSTED: 2026-10-09
% ATTEMPT_STATUS: unresolved
% ENTRY_KIND: research_attempt
\documentclass[11pt]{article}
\usepackage[margin=1in]{geometry}
\usepackage{amsmath,amssymb,amsthm,hyperref}
\newtheorem{proposition}{Proposition}
\newtheorem{lemma}{Lemma}
\title{Building effective state capacity: a research attempt}
\author{GPT-6 (Codex)}
\date{9 October 2026}
\begin{document}
\maketitle
\noindent\textbf{Status: unresolved.} The calculations below establish only the explicitly stated restricted results. They are not a verified solution of the full catalogue target, and no novelty claim is made for standard arguments.

\section{Target and a tractable boundary of the problem}
The canonical target chooses capability investments to maximize discounted public-service benefits while bounding predation and respecting political implementation. It also asks which gains survive withdrawal of support or political turnover. A tax-revenue increase does not by itself measure the three capability coordinates. The attempt first solves a finite, fully observed control version and then studies durability and identification. It does not assume that a politician obeys an unconstrained social planner.

Fix horizon $H$, finite states $s=(k,z,p)$, initial distribution $\nu$, discount $\beta\in(0,1)$ and finite menus $A_t(s)$. An action belongs to this menu only if its resource cost is affordable and its political implementation has already been established in a specified incentive model. Let $P_t(s'\mid s,a)$ be a known transition law, $c_t(s,a)$ net real benefit and $r_t(s,a)\geq0$ predation loss. All are bounded. Public randomization over feasible policies is permitted. These additional assumptions are strong, but make the optimization question exact.

\section{Occupation measures give an exact finite solution}
Set $x_t(s,a)=\Pr(S_t=s,A_t=a)$. Feasible occupation measures satisfy
\[
 x_t(s,a)\geq0,\quad \sum_a x_0(s,a)=\nu(s),\quad
 \sum_a x_{t+1}(s',a)=\sum_{s,a}P_t(s'\mid s,a)x_t(s,a).
\]
The constrained problem becomes the finite linear program
\[
 \max_x\sum_{t,s,a}\beta^t c_t(s,a)x_t(s,a),\qquad
 \sum_{t,s,a}\beta^t r_t(s,a)x_t(s,a)\leq\bar R.
\]
An empty feasible set is reported as such. Otherwise the flow polytope is bounded and closed, since every time slice has total mass one. A maximizer exists.

\begin{proposition}
Every feasible history-dependent randomized policy induces a feasible $x$. Conversely every feasible $x$ is implemented by the Markov randomized rule
$\pi_t(a\mid s)=x_t(s,a)/\sum_bx_t(s,b)$ on states of positive mass, with arbitrary feasible choices elsewhere.
\end{proposition}
\begin{proof}
The first statement is the law of total probability. For the converse the initial marginal agrees with $x_0$; applying the displayed flow equation inductively proves equality of all subsequent state-action marginals. Both objectives depend only on these marginals.
\end{proof}
Therefore memory beyond the state does not improve the objective in this benchmark. This is a standard control reduction, not a solution of endogenous political incentives.

The dual has a useful economic interpretation. For $\lambda\geq0$, solve
\[
 V_t^\lambda(s)=\max_{a\in A_t(s)}\{c_t(s,a)-\lambda r_t(s,a)
       +\beta\sum_{s'}P_t(s'\mid s,a)V_{t+1}^\lambda(s')\},
 \quad V_{H+1}^\lambda=0.
\]
Linear-program duality gives the constrained value as
$\inf_{\lambda\geq0}\{\lambda\bar R+\sum_s\nu(s)V_0^\lambda(s)\}$.
The multiplier is a shadow price per unit of discounted predation, not a claim that predation and service benefits are ethically interchangeable outside the announced objective.

With one resource inequality, some optimal policy can be implemented by randomizing initially between at most two deterministic history policies. Indeed, the finite set of deterministic policies generates a finite set of pairs $(R,C)$; randomization produces its convex hull in the plane. The highest feasible point lies at a vertex or on an edge. At most the two endpoints are needed. This does not mean independently mixing actions from two arbitrary policies preserves political incentive constraints: the initial lottery must itself be implementable.

\section{A transparent durability calculation}
Let a unit investment at date zero produce capability $k_t=d^t$, $0\leq d\leq1$, after support ends. Political control follows a row-stochastic matrix $P$, initially $p_0$, and service benefits and predation per unit capability are vectors $v$ and $r\geq0$. Assume retention is independent of political control; state-dependent destruction needs a different matrix. The discounted effects are
\[
 C_H=\sum_{t=0}^H(\beta d)^t e_{p_0}^{\mathsf T}P^tv-c_0,
 \qquad R_H=\sum_{t=0}^H(\beta d)^t e_{p_0}^{\mathsf T}P^tr.
\]
For an infinite horizon they become $e_{p_0}^{\mathsf T}(I-\beta dP)^{-1}v-c_0$ and the corresponding expression for $r$. The Neumann series converges because $\beta d<1$ and $P$ has spectral radius one. This proves that technical durability alone is insufficient: even $d=1$ may amplify predation when turnover places substantial mass on coercive regimes. Conversely rapid turnover does not mechanically destroy benefits when both regimes use the capability well.

For example, let an investment increase capability by $u\geq0$, cost $au$, yield marginal service value $B$, and predation $D>0$ using the preceding formulas. Under a budget $u\leq U$, the optimal amount is zero if $B\leq a$ and otherwise $\min\{U,\bar R/D\}$. A more durable technology can have a smaller feasible scale if its predation multiplier grows faster. This elementary example shows why maximizing persistence or collection alone need not solve the constrained objective.

\section{Why observed collection does not identify capability}
Suppose an administrative experiment measures only revenue $T=\tau k$, where $\tau$ is the effective extraction rate and $k$ the measured tax base times enforcement capability. Baseline $(\tau,k)=(1,1)$ gives $T=1$. A randomized tool can produce $T=2$ either through $(1,2)$, a capability improvement, or $(2,1)$, an extraction increase. Choosing identical revenue noise makes the entire observed trial law identical. The second model can have larger coercive loss and no durable capability gain. Assignment randomization identifies the revenue effect in both models, not which mechanism produced it.

Direct audit accuracy, case-resolution times, independent coercion measures, and post-turnover follow-up would eliminate portions of this equivalence class. For randomization with spillovers, cluster treatment or a justified exposure map is needed; a unit-level difference alone does not identify national scaling. If transition laws are only bounded, the occupation program for one chosen transition is not a robust solution. Nature's timing, whether ambiguity is rectangular, and permissible learning must be stated before dynamic robust optimization is valid.

\section{Assessment and primary sources}
The finite known-model optimization, two-policy mixture result, and turnover multiplier are proved above. The missing original-scope steps are political implementability of menus, identification of transitions and predation, and extrapolation across environments. These cannot be inferred from the algebra.

Timothy Besley and Torsten Persson (2009), \emph{The Origins of State Capacity: Property Rights, Taxation, and Politics}, AER 99(4), 1218--1244; primary publisher abstract inspected for the dependence of feasible policy on past investment:
\url{https://www.aeaweb.org/articles?id=10.1257/aer.99.4.1218}.
Their \emph{State Capacity, Conflict and Development}, NBER Working Paper 15088 (2009 working-paper version), primary abstract inspected for the political-conflict research questions:
\url{https://www.nber.org/papers/w15088}. These sources motivate the model; they are not cited as proofs of the particular constrained program here.

\end{document}
